\chapter{原子能发电厂动力设备及运行}

\section{思考题}

\begin{problem}
中子轰击重原子核时，受激核可能发生哪些核物观变化？迄今为止，世界各国大多采用 \(^{235}_{92}U\) 做核燃料的主要原因是什么？
\end{problem}

\begin{solution}
中子轰击重原子核时，受激核可能发生重原子核裂变现象，其内部的结合能（核力）会大量释放出来。

迄今为止，世界各国大多采用 \(^{235}_{92}U\) 做核燃料的主要原因是：
\begin{enumerate}
  \item 它是自然界存在的易裂变元素。
  \item 只有 \(^{235}U\) 的临界裂变能较小，且天然存在，非常适合作为反应堆靶核的核燃料元素。
\end{enumerate}
\end{solution}

\begin{problem}
从原子反应式 \(^{235}_{92}U + ^{1}_{0}n \rightarrow ^{137}_{56}Ba + Kr^{97}_{36} + 2 \times ^{1}_{0}n\)，说明核裂变释放巨大核能的基本原理。
\end{problem}

\begin{solution}
核裂变释放巨大核能的基本原理是：
\begin{enumerate}
  \item 质子和中子组成的稳定原子核的质量总是小于分离的中子和质子的质量总和，即在裂变重组过程中发生了质量亏损。
  \item 亏损掉的质量，以结合能的形式释放了出来。
  \item 根据爱因斯坦能量方程式 \(E=\Delta mc^{2}\)，这些质量亏损主要通过裂变碎片的动能形式转化为了巨大的热能。
\end{enumerate}
\end{solution}

\begin{problem}
何谓"中子俘获"、"链式裂变反应"、"中子慢化"？核裂变反应堆中，慢化剂应具备哪些基本特点？
\end{problem}

\begin{solution}
\begin{enumerate}
  \item \textbf{中子俘获}：由于中子不带电，能很容易地进入原子核中去，并与原子核结合，成为击破重原子核最为理想的粒子，这一过程体现了原子核对中子的俘获。
  \item \textbf{链式裂变反应}：易裂变的靶核（如 \(^{235}U\)）在核裂变开始后，便能像链条一样自动地、一环扣一环地持续进行下去的核裂变反应，称为自持链式裂变反应。
  \item \textbf{中子慢化}：利用慢化剂使能量大的快中子，消耗中子的动能，降低中子的速度到热中子状态的过程。
\end{enumerate}

慢化剂应具备的基本特点：
\begin{enumerate}
  \item 原子核质量轻（因为作为慢化剂的核靶原子量越小，中子每次与它散射时损失的能量就越大）。
  \item 要求慢化剂的密度大。
\end{enumerate}
\end{solution}

\begin{problem}
什么是自持链式裂变反应的"反应性"和"临界"？如何调节核裂变反应堆反应性的高低、从而改变堆输出功率的大小？
\end{problem}

\begin{solution}
\begin{enumerate}
  \item \textbf{反应性}：在反应堆的控制技术上，常用反应性 \(\rho\) 作为参量来量度堆内中子的增减，公式为 \(\rho=(K-1)/K\)（其中 K 为中子增殖系数）。
  \item \textbf{临界}：当 \(K=1\) 且反应性 \(\rho=0\) 时，新生一代中子与原来一代中子数目相等，链式反应刚好能自行维持下去，裂变反应的规模既不扩大也不减小，保持原状，称这种情况为反应堆的临界状态。
\end{enumerate}

调节方法：
\begin{enumerate}
  \item 主要通过移动控制棒（由易吸收中子的材料如镉或碳化硼制成）来调节中子的数量。增加或减小堆内的反应性 \(\rho\)，就能改变原子核裂变反应的规模，从而达到调节反应堆输出功率的目的。
  \item 此外，近代核反应堆通常会在冷却剂中适当加入硼酸，运行中只要改变冷却剂中硼的浓度，也可调整堆内反应性的高低。
\end{enumerate}
\end{solution}

\begin{problem}
什么是"裂变中子倍增时间"？如何使核裂变反应堆内裂变中子的倍增时间延长？为什么要延长裂变中子的倍增时间？
\end{problem}

\begin{solution}
\begin{enumerate}
  \item \textbf{裂变中子倍增时间}：反应堆内中子数增加 1 倍所需要的时间。
  \item \textbf{如何延长}：在设计反应堆时，使堆内中子剩余增殖系数控制在 \(0<(K-1)<0.7\%\)。这样堆内的核反应就需要靠瞬发中子加上寿命较长的缓发中子共同作用才能发展下去，从而使堆内中子的倍增时间延长达几秒钟以上。
  \item \textbf{为什么要延长}：如果仅靠瞬发快中子，其倍增时间常在几万分之一秒内，机械运动的控制棒无论如何是对付不了的，会导致反应堆处于瞬发超临界态的失控危险状态。延长倍增时间，才能确保能够通过控制棒的机械移动来安全控制堆内裂变反应和堆功率。
\end{enumerate}
\end{solution}

\begin{problem}
以轻水为冷却剂和慢化剂的压水式核反应堆核电厂，其核岛部分由哪些主要设备组成？简述一回路系统的原则性布置和循环过程。
\end{problem}

\begin{solution}
\begin{enumerate}
  \item \textbf{核岛主要设备}：由压水堆本体和一回路系统设备组成，主要包括反应堆本体、蒸汽发生器、一回路循环泵及其附属设备、连接管路，以及稳压器等。
  \item \textbf{一回路系统原则性布置与循环过程}：以反应堆为核心，核燃料在反应堆中进行可控链式裂变反应；冷却剂将裂变产生的大量热量带出反应堆，再通过蒸汽发生器将热量传给水，水被加热成蒸汽供汽轮机拖动发电机发电；冷却剂释热后，通过冷却剂循环主泵送回反应堆去继续吸热，从而不断地将核裂变释放的热能引导出来；整个一回路的压力则依靠稳压器来维持稳定。
\end{enumerate}
\end{solution}

\begin{problem}
与火力发电厂相比，压水堆核电厂的常规蒸汽发电系统的工质参数、设备结构、系统布置等，有哪些主要不同？
\end{problem}

\begin{solution}
\begin{enumerate}
  \item \textbf{工质参数}：新汽为饱和蒸汽或微湿蒸汽，工质为低压饱和蒸汽，因此工质流量和容积流量都比火电厂大得多。
  \item \textbf{设备结构}：必须先行汽水分离；高压缸排蒸汽湿度较大，必须再加热，其再热器的特点是用新蒸汽或抽汽作为加热蒸汽；在湿蒸汽区域的各级叶栅，需采取相应的去湿措施；在相同功率下，核电机组比热电机组的体积和质量更大、效率更低，且必须设置中间汽水分离器。
  \item \textbf{系统布置}：需有蒸汽母管；需用整机旁路系统；给水系统为母管式单元制系统；须有连续放射性监测装置；且核电厂汽轮机组的循环冷却水需求量大。
\end{enumerate}
\end{solution}

\begin{problem}
什么是放射性？激发态及裂变过程中的核素，放射出的 \(\alpha\)、\(\beta\)、\(\gamma\) 射线是由何种粒子组成的？它们各具有怎样的辐照特点？
\end{problem}

\begin{solution}
\begin{enumerate}
  \item \textbf{放射性}：是指放射性同位素在衰变、或处于不稳定的高能量状态时发生转变，释放出射线（如粒子流或电磁波），并可能转变为新原子核的特性。
  \item \textbf{\(\alpha\) 射线}：是由两个质子及两个中子组成的氦原子核组成，带正电荷。辐照特点是穿透物质的本领比 \(\beta\) 射线弱得多，容易被薄层物质阻挡，但是它有很强的电离作用。
  \item \textbf{\(\beta\) 射线}：由核内中子过多时释放出的电子流（\(\beta^-\)衰变）或核内质子过多时释放出的正电子流（\(\beta^+\)衰变）组成。辐照特点是穿透本领比 \(\alpha\) 射线强，但在空气中容易被吸收，通常一层薄的有机玻璃就能有效地阻止 \(\beta\) 射线的外照射。
  \item \textbf{\(\gamma\) 射线}：是一种波长极短的电磁波，不带电且没有静止质量。辐照特点是具有极强的穿透本领，过量的 \(\gamma\) 辐照会对人体和设备带来很大的伤害。
\end{enumerate}
\end{solution}

\begin{problem}
核电厂设置有几道安全屏蔽防护？它们各是如何防止放射性物质外逸的？从事核电设备运行、维护和检修的工作人员，应具有哪些基本的防辐照知识？
\end{problem}

\begin{solution}
核电厂设置有四道安全屏蔽防护：
\begin{enumerate}
  \item \textbf{第一道：核燃料本身}。制成物理、化学性能十分稳定的二氧化铀陶瓷块，能把裂变产物的 98\% 以上保持在芯块内，只要不熔化，裂变产物就不会大量跑出。
  \item \textbf{第二道：燃料元件的包壳}。用优质锆合金材料制成，能将逸出的少量裂变产物基本上保持在包壳密封之内。
  \item \textbf{第三道：压力容器}。把燃料组件等完全封闭起来，冷却剂循环时不与核燃料直接接触，即使包壳少量破漏，放射性物质也只会扩散到封闭的一回路中。
  \item \textbf{第四道：安全壳}。庞大的安全壳把整个一回路的设备系统包覆起来，即使一回路出现破裂或渗漏，放射性物质也不会逸出到外部环境中去。
\end{enumerate}

工作人员应具备的基本防辐照知识：
\begin{enumerate}
  \item \textbf{外部防护}：在有有效屏蔽设施的情况下，必须具备熟练的技术以缩短受照条件下的操作时间，并尽量采用遥控和远距离操作；在辐射区内的连续工作时间不允许超过容许照射剂量率规定的时间。
  \item \textbf{内部防护}：所有受照操作时必须配备有效的防护用具，严禁在工作场所吸烟、饮水和进食，并依靠精确的仪器对禁区范围进行连续、严密的放射性污染监测。
\end{enumerate}
\end{solution}

\begin{problem}
什么是核电厂的"三废"？核电厂的三废物质和乏燃料的基本处理方法有哪些？
\end{problem}

\begin{solution}
\textbf{"三废"}：核电站在运行时，不可避免地会产生的带有放射性的废水、废气和固态废料。

处理方法：
\begin{enumerate}
  \item \textbf{废水}：主要采用储存衰变法（使短半衰期放射性同位素完成衰变）、蒸发法（浓缩后加入水泥等凝固剂固化封装）以及离子交换法（用树脂吸附放射性物质后固化处理）。
  \item \textbf{废气}：采用衰变法（加压后送入衰变箱经过 60-100 天衰变后稀释排入大气）以及过滤吸附法（用玻璃纤维滤纸和颗粒活性碳层去除气溶胶微粒和碘蒸气）。
  \item \textbf{固体废物}：可燃质低放射固形物焚烧后，与松散固态物质一并压缩装桶，与其他密实废物一起深埋于地下，或沉入深海；低、中放射性核废物可存放入近地表处置场。
  \item \textbf{乏燃料}：送入专门工厂进行"后处理"，分离和回收其中有利用价值的同位素（如未完全燃尽的 U-235 和新生成的 Pu-239）；剩余的带有放射性的废物经浓集后，像固体废物一样作永久性处理（或继续贮存）。
\end{enumerate}
\end{solution}
